\documentclass[10pt,a4paper]{amsart}
\usepackage[margin=19mm]{geometry}
\usepackage{amsmath,amssymb,mathtools}
\usepackage[scr=rsfs,cal=euler]{mathalfa}
\usepackage{xfrac}
\usepackage[unicode,hidelinks,bookmarks=false]{hyperref}
\newcommand{\ie}{i.e.\ }
\newcommand{\cf}{cf.\ }
\newcommand{\resp}{respectively\ }
\newcommand{\Subsection}{\subsection}
\providecommand{\nobreakdash}{\nobreak\hbox{-}}
\setlength{\emergencystretch}{2em}
\multlinegap0pt
\usepackage{amssymb}
\usepackage{mathtools}
\newtheorem{lemma}{Lemma}
\usepackage{cleveref}
\crefname{lemma}{Lemma}{Lemmas}
\newcommand{\IM}{\Delta}
\title{A greedy open-orbit criterion for solvable algebraic group actions, with applications to Lusztig's nilpotent varieties}
\author{Erez Lapid}
\date{Source: arXiv:2608.14413v1 (CC BY 4.0)}
\begin{document}
\maketitle

\noindent\emph{Source and licence.} A greedy open-orbit criterion for solvable algebraic group actions, with applications to Lusztig's nilpotent varieties. Version arXiv:2608.14413v1 (CC BY 4.0), distributed by arXiv under the Creative Commons Attribution 4.0 International licence, http://creativecommons.org/licenses/by/4.0/, as recorded in the arXiv licence metadata for this exact version and shown on the abstract page https://arxiv.org/abs/2608.14413v1. Full licence notice: \texttt{licenses/arxiv-2608.14413-CC-BY-4.0.txt}.\par\smallskip

\noindent\textit{Editorial scope.} Counted unit: the article's lemma minviol (source0.tex lines 1684--1699). The article's complete original proof of it is delivered with it (source0.tex lines 1710--1824, one \texttt{proof} environment of 115 lines, which the article places away from the statement). No mathematics is written, repaired or re-derived: the statement and the proof are the article's own lines, in the article's own notation, and no retained line is substituted or re-flowed.  Retained uncounted as the source-local context the proof needs: lines 1677--1679.  Nothing was omitted from any retained span, so the package carries no omission record.  No retained line contains a citation, so the wrapper carries no bibliography and no bibliography entry.  No retained line contains a figure environment, an image-inclusion command or drawing code, so no image is delivered and the package declares no asset dependency.  Editorial formatting only, in addition: an AMS \texttt{amsart} preamble that restates the article's own declarations and macros used by the retained lines, namely the environments \texttt{lemma} on separate counters, together with the standard packages those lines need.  The retained environments print in this document as Lemma 1 in order of appearance, whereas the article prints the same items with the numbering of its own full document.  The complete native source of the article as distributed in the arXiv e-print for this exact version is bundled unchanged as \texttt{sources/207668/source0.tex}.
\begin{equation} \label{eq: famI}
(k,\ D_1,\ D_2,\ \dots,\ D_t,\ 1)\in S_k,\qquad t\ge1,
\end{equation}

\begin{lemma} \label{lem: minviol}
The minimal violators are precisely the following permutations.
\begin{enumerate}
\item \label{part: mvI} The thickenings, with $k_1=k_m=1$, of $(m,2,\dots,m-1,1)$, $m\ge4$, i.e., the permutations \cref{eq: famI} with $t\ge2$ runs.
\item \label{part: mvII} The permutations
\[
\tau_{r,s}:=(r+1,\,r-1,r-2,\dots,2,\,k,\,1,\,k-1,k-2,\dots,r+2,\,r)\in S_k,\quad k=r+s,\ r,s\ge2
\]
(for $r=s=2$ this is $3412$).
\item \label{part: mvIII} The permutations
\[
\upsilon_r:=(r+1,\ r+2,\ r,r-1,\dots,3,\ 1,\ 2)\in S_{r+2},\qquad r\ge3,
\]
i.e., the thickenings of $45312=\upsilon_3$ in which only the middle entry is thickened.
\end{enumerate}
\end{lemma}

\begin{proof}[Proof of \cref{lem: minviol}]
The proof is routine, if tedious.

It is convenient to argue in terms of the matrix $\IM=\IM_\sigma$, whose ones are $(i,\pi(i))$, $i=1,\dots,n$,
where $\pi=w_0\sigma$; thus $\sigma$ is smooth if and only if $\IM$ contains no occurrence of (the permutation
matrices of) $1324$ or $2143$, and in the list of the lemma the three families read, in terms of $\pi$:
\begin{align}
&\pi=(1,\ B_1,\dots,B_t,\ n),\label{eq: famIpi}\\
&\pi=(s,\ s+2,s+3,\dots,s+r-1,\ 1,\ n,\ 2,3,\dots,s-1,\ s+1),\label{eq: famIIpi}\\
&\pi=(2,\ 1,\ 3,4,\dots,n-2,\ n,\ n-1),\label{eq: famIIIpi}
\end{align}
where in \cref{eq: famIpi} the middle part is the partition of $\{2,\dots,n-1\}$ into $t$ \emph{increasing} runs of
consecutive integers arranged in decreasing order of their values (the three equations corresponding to the three
families of the lemma, in order).
Call the ones of $\IM$ in the first row, the last row, the first column and the last column the \emph{border} ones;
they are $f=(1,u)$, $l=(n,v)$, $c=(a,1)$ and $d=(b,n)$, where $u=\pi(1)$, $v=\pi(n)$, $a=\pi^{-1}(1)$, $b=\pi^{-1}(n)$.
Since deleting a border one (with its row and column) preserves any occurrence not containing it, and destroys any
occurrence containing it, we have:
\emph{$\sigma$ is a minimal violator if and only if $\IM$ contains an occurrence of $1324$ or $2143$, and every such
occurrence contains all four border ones.}

\emph{Step 1: coincidences among the border ones.}
The possible coincidences are $f=c$ ($u=1$), $f=d$ ($u=n$), $l=c$ ($v=1$), $l=d$ ($v=n$).
If $u=n$, then every occurrence contains $f$, i.e., its first letter is its largest one; but the patterns $1324$ and
$2143$ begin with their first and second smallest letters, respectively. Hence there is no occurrence at all,
contradicting non-smoothness. Similarly $v=1$ is impossible (the patterns do not end with their smallest letter).
Next suppose $u=1$ but $v\ne n$. Then $f=c$ and the border ones are $f$, $d=(b,n)$, $l=(n,v)$ with $v<n$; an
occurrence contains all three, and since its first letter is the smallest one ($1$ at position $1$), its pattern
begins with its minimal letter, so it must be $1324$; but then its last letter is its largest one, whereas the letter
at the last position is $v<n$, and $n$ occurs inside. This is a contradiction; hence $u=1$ implies $v=n$, and
symmetrically $v=n$ implies $u=1$. We are left with two cases: either $u=1$ and $v=n$, or $1<u,v<n$ and the four
border ones are distinct.

\emph{Step 2: the case $u=1$, $v=n$.}
Here $f=c=(1,1)$ and $l=d=(n,n)$. Let $W$ be the word $\pi(2),\dots,\pi(n-1)$.
For interior positions $i<j$, the quadruple $(1,i,j,n)$ is an occurrence (necessarily of $1324$) if and only if
$\pi(i)>\pi(j)$; such occurrences contain all border ones. Any other potential occurrence omits a border one and is
therefore forbidden:
\begin{itemize}
\item $(i,j,k,n)$, $1<i<j<k<n$: an occurrence of $1324$ if and only if $(\pi(i),\pi(j),\pi(k))$ forms a
$132$-pattern (an occurrence of $2143$ is impossible since the largest letter of $2143$ is its third one).
Hence $W$ must avoid $132$.
\item $(1,i,j,k)$: an occurrence of $1324$ if and only if $(\pi(i),\pi(j),\pi(k))$ forms a $213$-pattern
($2143$ is impossible since the second letter of $2143$ is its smallest one). Hence $W$ must avoid $213$.
\item occurrences within the interior contain a $132$- or $213$-pattern and are excluded by the above.
\end{itemize}
Thus, the case at hand holds if and only if $W$ avoids $132$ and $213$ and is not increasing.
We claim that $W$ avoids $132$ and $213$ if and only if $W$ is a concatenation of increasing runs of consecutive
integers, arranged in decreasing order of their values (so that non-increasing means $t\ge2$ runs, giving exactly the
family \cref{eq: famIpi}). Indeed, if $W$ has this form, then any ascent of $W$ lies within a single run, and both
patterns require an ascent $(x,y)$ together with a letter strictly between $x$ and $y$ located outside the positions
between them, which is impossible for a run of consecutive integers. Conversely, suppose $W$ avoids both patterns and
let $d$ be the position of the largest letter $t_0$ of $W$. Avoidance of $132$ for the triples $(i,d,k)$, $i<d<k$,
forces every letter after $d$ to be smaller than every letter before $d$; avoidance of $213$ for the triples
$(i,i',d)$ forces the prefix up to $d$ to be increasing. Hence this prefix consists of the largest letters in
increasing order, i.e., it is an increasing run of consecutive integers ending with $t_0$, and we conclude by
induction on the length applied to the remaining suffix.

\emph{Step 3: the case of distinct border ones.}
Every occurrence contains the four distinct border ones, hence consists exactly of them. The positions of $f,c,d,l$
are $(1,a,b,n)$ or $(1,b,a,n)$ in increasing order. If $b<a$ then the letter $n$ precedes the letter $1$ at the
second and third positions of the occurrence; but in $1324$ and $2143$ the largest letter does not precede the
smallest one. So there is no occurrence, a contradiction; hence $a<b$. The unique candidate occurrence
$(1,a,b,n)$ has the letters $(u,1,n,v)$, which form an occurrence (necessarily of $2143$) if and only if $u<v$.
Non-smoothness thus forces
\[
1<u<v<n,\qquad 1<a<b<n,
\]
and the requirement is that \emph{no other} quadruple is an occurrence. Let
$P$, $Q$, $R$ be the sets of non-border positions $i$ with $i<a$, $a<i<b$, $b<i$, respectively.
Forbidding the following quadruples (none of which contains all four border ones) yields, in turn:
\begin{itemize}
\item $(i,a,b,n)$, $i\in P$: letters $(\pi(i),1,n,v)$, an occurrence of $2143$ if and only if $\pi(i)<v$.
Hence $\pi(i)>v$ on $P$.
\item $(1,a,b,j)$, $j\in R$: letters $(u,1,n,\pi(j))$, an occurrence of $2143$ if and only if $\pi(j)>u$.
Hence $\pi(j)<u$ on $R$.
\item $(1,a,j,n)$, $j\in Q\cup R$: letters $(u,1,\pi(j),v)$, an occurrence of $2143$ if and only if $\pi(j)>v$.
Hence $\pi(j)<v$ for $j>a$; together with the previous item, $u<\pi(j)<v$ on $Q$.
\item $(1,i,j,b)$, $i<j$ in $P$: letters $(u,\pi(i),\pi(j),n)$, an occurrence of $1324$ if and only if
$u<\pi(j)<\pi(i)$. Since $\pi(i),\pi(j)>v>u$, this forces $\pi$ to be increasing on $P$.
Similarly, $(1,i,j,n)$ for $i<j$ in $Q$ and $(a,i,j,n)$ for $i<j$ in $R$ force $\pi$ to be increasing on $Q$ and
on $R$.
\item $(1,i,j,b)$, $i\in P$, $j\in Q$: letters $(u,\pi(i),\pi(j),n)$ with $u<\pi(j)<\pi(i)<n$ --- always an
occurrence of $1324$. Hence $P$ and $Q$ cannot both be nonempty. Similarly, $(a,i,j,n)$ for $i\in Q$, $j\in R$
has letters $(1,\pi(i),\pi(j),v)$ with $1<\pi(j)<\pi(i)<v$, so $Q$ and $R$ cannot both be nonempty.
\end{itemize}
Counting values yields $\#P=n-1-v$, $\#Q=v-u-1$, $\#R=u-2$.
If $Q=\emptyset$ then $v=u+1$, $b=a+1$, and the constraints above determine $\pi$ completely:
$\pi$ is given by \cref{eq: famIIpi} with $s=u$ and $r=n-u$ (the case $P=\emptyset$, resp.\ $R=\emptyset$,
corresponding to $r=2$, resp.\ $s=2$). If $Q\ne\emptyset$ then $P=R=\emptyset$, so $u=2$, $v=n-1$, $a=2$, $b=n-1$,
and $\pi$ is given by \cref{eq: famIIIpi}.

\emph{Step 4: sufficiency.}
In Step 2 it was already shown that the permutations \cref{eq: famIpi} with $t\ge2$ are minimal violators.
For \cref{eq: famIIIpi}, the only inversions of $\pi$ are at the positions $(1,2)$ and $(n-1,n)$; an occurrence of
$1324$ uses an inversion with a letter before it and a letter after it, which is impossible here, while an occurrence
of $2143$ uses two disjoint inversions, so the unique occurrence is $(1,2,n-1,n)$, which consists exactly of the four
border ones. For \cref{eq: famIIpi}, write $\pi=(s,\ P,\ 1,\ n,\ R,\ s+1)$, where $P$ (occupying the positions $2,\dots,a-1$,
$a=r$) and $R$ (occupying the positions $b+1,\dots,n-1$, $b=r+1$) are the increasing words with letters
$s+2,\dots,s+r-1$ and $2,\dots,s-1$, respectively. Abusively denoting positions by the blocks they belong to,
the inversions of $\pi$ are: $(x,a)$ with $x\in\{1\}\cup P$; $(x,y)$ with $x\in P$, $y\in R\cup\{n\}$;
$(1,y)$ with $y\in R$; and $(b,y)$ with $y\in R\cup\{n\}$.
An occurrence of $2143$ consists of two disjoint inversions $(i_1,i_2)$, $(i_3,i_4)$ with $i_2<i_3$ and
$\pi(i_2)<\pi(i_1)<\pi(i_4)<\pi(i_3)$. From the catalogue, an inversion $(i_1,i_2)$ with $i_2<i_3\le b$ must have
$i_2=a$, so $\pi(i_2)=1$ and $\pi(i_1)\in\{s\}\cup P$; and then $i_3=b$ ($i_3=a$ is excluded, and there is no
inversion starting in $R$). If $i_4\in R$ then $\pi(i_4)\le s-1<s\le\pi(i_1)$, a contradiction; so $i_4=n$ and
$\pi(i_4)=s+1$, forcing $\pi(i_1)=s$, i.e., $i_1=1$. Thus the unique occurrence of $2143$ is $(1,a,b,n)$.
An occurrence of $1324$ consists of an inversion $(i_2,i_3)$ together with positions $i_1<i_2$, $i_4>i_3$ such that
$\pi(i_1)<\pi(i_3)$ and $\pi(i_4)>\pi(i_2)$. Running over the catalogue: if $i_3=a$ then $\pi(i_1)<1$, which is
impossible; if $i_2\in\{b\}$ then $\pi(i_4)>n$, impossible; if $i_2\in P$ and $i_3\in R\cup\{n\}$ then either
$i_3=n$ is the last position (no room for $i_4$) or $\pi(i_4)>\pi(i_2)\ge s+2$ with $i_4$ after a position of $R$,
whereas all letters there are at most $s+1$, impossible; and if $i_2=1$, $i_3\in R$, then there is no room for
$i_1$. Hence there is no occurrence of $1324$. Since every occurrence contains all four border
ones, all four deletions are smooth, while $\sigma$ itself is not.
\end{proof}

\end{document}
